\chapter{Kết luận và hướng phát triển}
\label{ch:conclusion}

\section{Trả lời các câu hỏi nghiên cứu}

\paragraph{RQ1.} Với 256 mẫu có nhãn mỗi tác vụ, các bộ định tuyến nâng điểm trung bình từ 57.4 lên $69.9\pm0.1$ và mức tăng so với lớp đa số từ 18.4\% lên 59.1\%. Khoảng cách rõ rệt với các chuyên gia (78.2) vẫn còn.

\paragraph{RQ2.} Điểm trung bình tăng từ 60.8 ở $k=16$ lên 64.6 ở $k=64$ và 69.9 ở $k=256$, và chưa chững lại. Ngân sách nhỏ không ổn định giữa các seed.

\paragraph{RQ3.} Các hệ số cố định chỉ đạt 58.4, nên riêng các hệ số tĩnh tốt hơn không giải thích được cải thiện, dù ablation chưa tách được tính phụ thuộc đầu vào khỏi sức chứa.

\paragraph{RQ4.} Nâng cấp full MLP và intermediate-only đạt điểm trung bình tương tự với các thất bại khác nhau, lần lượt ở RTE và MRPC. Intermediate-only cần 47.2\% dung lượng của năm mô hình, so với 73.1\%. Ở batch 32, intermediate-only và full MLP chậm hơn một BERT khoảng 2.9 và 4.5 lần.

\section{Đánh giá các mục tiêu}

Bảng~\ref{tab:goals} đánh giá các mục tiêu ở Mục~\ref{sec:problem} dựa trên bằng chứng.

\begin{table}[htbp]
  \centering\tablestyle
  \caption{Các mục tiêu của bài toán so với kết quả.}
  \label{tab:goals}
  \begin{tabularx}{\linewidth}{>{\raggedright\arraybackslash}p{3.2cm} X >{\raggedright\arraybackslash}p{5.2cm}}
    \toprule
    Mục tiêu & Bằng chứng & Đánh giá\\
    \midrule
    O1. Tích hợp động & Trọng số MLP theo từng đầu vào; hệ số cố định đạt 58.4, so với 69.9 & Đạt theo thiết kế; đóng góp của tính phụ thuộc đầu vào được ủng hộ nhưng chưa được tách riêng\\
    O2. Gần các mô hình riêng theo tác vụ & 69.9 so với 78.2, mức tăng 59.1\%; RTE và MRPC vẫn yếu; không tác vụ nào đạt mức chuyên gia & Đạt một phần\\
    O3. Ít tham số lưu trữ hơn & 73.1\% (full MLP) hoặc 47.2\% (intermediate-only) của năm mô hình & Đạt, đổi lại độ trễ cao hơn\\
    \bottomrule
  \end{tabularx}
\end{table}

\section{Đóng góp chính và ý nghĩa thực tiễn}

Đóng góp mang tính thực tiễn, và không có thuật toán mới nào được đề xuất. Nghiên cứu này chuyển WEMoE sang encoder BERT chuẩn hóa sau với định tuyến có tính đến mặt nạ attention, áp dụng quy trình tách biệt dữ liệu của từng giai đoạn, và đánh giá hai phạm vi nâng cấp, ba ngân sách nhãn và một ablation không phụ thuộc đầu vào so với mốc lớp đa số. Trong thực tế, hợp nhất phụ thuộc đầu vào vượt rõ rệt hợp nhất tĩnh về trung bình nhưng không vượt các chuyên gia riêng theo tác vụ, và lý do dùng nó chủ yếu nằm ở ràng buộc lưu trữ.

\section{Hướng phát triển}

Hợp nhất bốn tác vụ không có RTE sẽ kiểm tra liệu RTE có can nhiễu các tác vụ khác không. Ngân sách lớn hơn ($k=512$ và $1024$) sẽ cho thấy điểm có chững lại không, và thêm bước FSFT sẽ kiểm tra biến thể intermediate-only. Đo chuẩn của véc-tơ tác vụ theo từng loại mô-đun sẽ kiểm chứng lý do nâng cấp MLP trên BERT. Một quy trình nghiêm ngặt sẽ giới hạn dev ở $k$, và thêm seed cho các chuyên gia và cách chia dữ liệu sẽ cho ước lượng đáng tin cậy hơn. Phân tích theo từng mẫu sẽ cho thấy biến thiên định tuyến trong một tác vụ đi theo yếu tố nào.
